% !PRINTED 67
\weberpage{67}
%<<CONT

\gesperrt{III. The permutations of the first kind are composed of an even number of transpositions, and those of the second kind of an odd number.}

The so-called identical permutation, which leaves \(\mathfrak A\) unchanged, is therefore also to be counted as of the first kind.

From this one obtains further the conclusion: in however many different ways one may derive \(\mathfrak A'\) from \(\mathfrak A\) by transpositions, the number of these transpositions is in all these ways consistently even or odd, according as \(\mathfrak A'\) belongs to the first or to the second kind.

If in all arrangements
\[
 \mathfrak A,\mathfrak A',\mathfrak A'',\ldots
\]
of the \(n\) elements we perform one transposition, say \((1,2)\), then each of these arrangements passes into a definite other one, say \(\mathfrak A\) into \(\mathfrak B\), \(\mathfrak A'\) into \(\mathfrak B'\), \(\mathfrak A''\) into \(\mathfrak B''\), and so on; and if the same transposition is repeated once more, then \(\mathfrak B\) again passes into \(\mathfrak A\), \(\mathfrak B'\) again into \(\mathfrak A'\), and so on. Hence the arrangements \(\mathfrak B,\mathfrak B',\mathfrak B'',\ldots\) are all distinct from one another, and consequently in their totality agree with the totality of the \(\mathfrak A\)'s. Since, as we have seen above, all difference-products
\[
 P,P',P'',\ldots
\]
which are formed from \(P\) with the different arrangements \(\mathfrak A,\mathfrak A',\mathfrak A'',\ldots\) change their sign under a transposition, it follows that to each \(\mathfrak A\) of the first kind there corresponds a \(\mathfrak B\) of the second kind, and to each \(\mathfrak A\) of the second kind a \(\mathfrak B\) of the first kind.

\gesperrt{IV. Accordingly the number of arrangements of the first kind is just as large as the number of arrangements of the second kind, namely} \({1\over2}\Pi(n)\)\footnote{These theorems have here been obtained from the consideration of the product \(P\), hence from a numerical quantity. How one can arrive at the same results without using such a function will be seen in the fourteenth section.}.

For \(n=3\) we have the following six arrangements, of which the first horizontal row forms the first kind:
\begin{equation}
\begin{array}{ccc}
(1,2,3),&(2,3,1),&(3,1,2),\\
(3,2,1),&(2,1,3),&(1,3,2).
\end{array}
\tag{3}
\end{equation}
